\chapter{Introduction}

\section{The long-read mapping problem}

Given a \emph{reference} sequence (or a set of reference segments, e.g.\ chromosomes/contigs)
$T = T_1, \dots, T_s$ over the alphabet $\Sigma = \{A, C, G, T\}$, and a \emph{query} read $P$
(typically $1$--$100+$\,kbp for long-read sequencing technologies such as PacBio HiFi or Oxford
Nanopore), the mapping problem is: find the location(s) $(i, [l, r))$ --- segment index $i$ and
half-open interval $[l, r) \subseteq [0, |T_i|)$ --- such that $T_i[l..r)$ is the region of the
reference that $P$ most likely originated from, together with a confidence score (mapping
quality, MAPQ).

This differs from \emph{alignment}: mapping only needs to find the correct region and strand, not
compute a base-level edit script (CIGAR). Mapping is the first, and usually the computational
bottleneck, phase of most long-read tools (minimap2, Winnowmap, MashMap, mapquik, and shmap all
separate mapping from an optional downstream alignment phase).

The naive solution --- align $P$ against every possible substring of $T$ with a full dynamic
program (edit distance / Smith--Waterman) --- costs $O(|T|\cdot|P|)$ per read and is many orders
of magnitude too slow for genome-scale references ($|T| \sim 10^9$--$10^{10}$) and
population-scale read sets ($10^5$--$10^7$ reads). All practical long-read mappers instead use a
\textbf{seed-and-extend} (or \emph{seed-chain-extend}) strategy: cheaply find a small number of
short, exact-matching anchors (\emph{seeds}) between $P$ and $T$, then use only the anchors'
positions --- not the sequence content --- to decide where $P$ came from.

\section{shmap's approach in one paragraph}

shmap represents both the reference and each read as a \textbf{FracMinHash sketch}: a
size-proportional, hash-thresholded subset of the sequence's $k$-mers (Chapter~\ref{ch:theory}
develops this precisely). Reference $k$-mers are indexed once into a hash table mapping
$k$-mer hash $\to$ list of reference positions. For each read, its own sketch is grouped into
\emph{seeds} (rarest $k$-mers first) and matched against the index; matches are accumulated into
overlapping, fixed-length \emph{buckets} along each reference segment. A cheap, provably-safe
\emph{seed heuristic} upper bound is used to prune buckets that cannot possibly beat the current
best candidate before any expensive refinement is attempted on them. Buckets that survive pruning
are refined with an exact two-pointer / sliding-window sweep (or, in alternate configurations, a
longest-common-subsequence chain) to produce a Jaccard or containment similarity score; the best
and second-best surviving mappings are used to compute a mapping-quality estimate; and the result
is emitted in PAF format with a rich set of custom diagnostic tags.

Two design choices set shmap apart from more familiar tools like minimap2:
\begin{enumerate}
  \item \textbf{FracMinHash sketching instead of minimizer windows.} Rather than picking one
    minimal-hash $k$-mer per sliding window (a \emph{minimizer} scheme), shmap keeps every k-mer
    whose hash falls below a fixed fraction of the hash space. This makes sketch size
    \emph{proportional to sequence length} with a tunable ratio $r$ (flag \code{-r}), rather than
    proportional to (length / window size); see \S\ref{sec:fracminhash-vs-minimizer}.
  \item \textbf{Score-bound-driven bucket pruning instead of chaining every anchor.} Instead of
    running a full co-linear chaining dynamic program over every seed hit (as minimap2 does),
    shmap only ever considers small, fixed-geometry \emph{buckets} of the reference, and prunes
    the overwhelming majority of them using a cheap closed-form upper bound on the score they
    could possibly achieve, before doing any real work on them. This is the ``SH'' (seed
    heuristic) in the tool's Map-SHmap name.
\end{enumerate}

\section{How to use this document}

The chapters are ordered so that a reader with no prior background in sketch-based mapping can
start at Chapter~\ref{ch:theory} and, by the time they reach the per-module chapters
(Chapters~\ref{ch:sketching}--\ref{ch:scoring}), already understand \emph{why} every formula and
data structure exists, not just \emph{what} it computes. Each per-module chapter gives:
\begin{itemize}[nosep]
  \item the exact input/output contract of every function (types, units, invariants assumed and
    guaranteed);
  \item a boxed pseudocode listing that is language-independent and directly implementable;
  \item excerpts of the original C++ (from the \texttt{shmap} repository, files under
    \texttt{src/}) and of an independent Rust re-implementation, side by side, wherever the exact
    bit-level or control-flow behavior matters;
  \item worked numeric examples for the trickiest arithmetic (bucket geometry, the FracMinHash
    rolling hash, the seed-heuristic bound);
  \item boxed notes on defects discovered by cross-checking the two implementations against each
    other and against the algorithm's own stated invariants --- several of which are not merely
    cosmetic (Chapter~\ref{ch:defects} collects all of them with full detail).
\end{itemize}

Chapter~\ref{ch:end-to-end} then reassembles every module into one top-to-bottom pseudocode
listing of the per-read mapping routine, which is the single most useful reference for a
from-scratch implementation. Chapter~\ref{ch:appendix} collects notation, a symbol table, and a
condensed algorithm index for quick lookup while implementing.
